\begin{annexitem}{annex:p03.op.1}{thm:p03-op-initial}{Initial segments, successors and carriers}
\begin{proof}
All auxiliary clauses below are finitely many first-order clauses.
Regularity supplies their conjunctions and the needed substitutions.
For infinite \(\kappa\), adding finitely many clauses preserves the
bound. For finite \(\kappa\geq1\), first conjoin the original
finite theory, and conjoin the auxiliary clauses with it; the result
is one sentence. Thus none of the constructions increases the
permitted number of sentences.

Suppose \(T\) pins \(\alpha\) on \(P,<\). Add a new unary
predicate \(J\), required to be a downward-closed subset of \(P\),
and designate \(J\) with the restricted order. Every resulting
carrier is an initial segment of a well-order. In a witness of
type \(\alpha\), choose its initial segment of any type
\(\xi\leq\alpha\), including the empty segment and the whole
carrier. The same expanded theory therefore pins every such
\(\xi\). This proves downward closure without assuming that
arbitrary subsets form substructures for the original vocabulary.

For successor, use \cref{thm:p03-general-fibres} to place a model
of \(T\) in the fibre at a distinguished constant \(d\).
Require \(\widehat\varphi(d)\) for \(\varphi\in T\).
Its carrier predicate and order are first-order definable in the
auxiliary structure by the displayed fibre interpretation; denote
them by \(P_d\) and \(<_d\). Add a point \(t\notin P_d\)
and designate the carrier \(P_d\cup\{t\}\), ordered by
\(<_d\) together with all pairs \((a,t)\), \(a\in P_d\).
A fresh unary and binary symbol can name these two definitions
using first-order biconditionals. Every model gives a successor
well-order. For a witness take a disjoint copy of the given
\(T\)-model as fibre and use an index point outside it as \(t\).
The fibre construction permits these additional ambient points
without changing its induced model. Its order type is
\(\alpha+1\). Finite carrier orders, including types zero and
one, also have direct first-order one-sentence descriptions.

A strict order on its field is a carrier order by defining the
carrier to be that field. Conversely, name the restriction of a
carrier order to its carrier by a fresh binary symbol. Its field
has the same order type unless the carrier is a singleton; in
that case the field is empty. Thus a theory pinning a carrier
type other than one pins that type under the field convention,
and every field-pinnable type is carrier-pinnable. No strict
order has a singleton field: any field element participates in
a strict pair with a different element. Since the finite
pinnable types are unbounded in \(\omega\), omitting the type
one changes neither supremum.

Successor closure implies that a bounded supremum \(W\) is
not attained. It is at least \(\omega\) and is a limit ordinal.
For each \(\xi<W\), the definition of supremum supplies a
pinnable ordinal at least \(\xi\), and downward closure pins
\(\xi\). No ordinal at least \(W\) can be pinned. These
observations also give the stated exact list for fields.
\end{proof}
\end{annexitem}

\begin{annexitem}{annex:p03.op.2}{thm:p03-op-threshold}{The cofinal-model threshold}
\begin{proof}
Suppose first that \(W_\kappa=W\) is bounded. Assume, towards
a contradiction, that a theory \(T\) of at most \(\kappa\)
sentences has only well-ordered carrier models, with types below
\(W\) cofinal in \(W\). Construct a new theory as follows.
Use \cref{thm:p03-general-fibres} with fresh index predicate
\(J\) and a strict linear order \(\prec\) on \(J\).
For each \(\varphi\in T\), require
\[
 \forall d\,(J(d)\longrightarrow\widehat\varphi(d)).
\]
Every index fibre is therefore a model of \(T\). Let \(P_d,<_{d}\)
be its definable carrier and order. Introduce a ternary relation
\(H(d,a,j)\). The first-order formula \(G(d)\) says that
\(H(d,-,-)\) is an order isomorphism from \(P_d,<_{d}\) onto
\(\{j\in J:j\prec d\},\prec\): require its pairs to lie in
these two carriers, exactly one image for each source point,
exactly one preimage for each target point, and preservation and
reflection of the two strict orders. These are finitely many
first-order clauses. Require also
\[
 \forall y\in J\ \exists d\in J\,
             (y\prec d\land G(d)).
\]
Together with the linear-order axioms this adds only finitely
many clauses. As in the preceding proof, finite conjunction
handles finite \(\kappa\); infinite \(\kappa\) is unchanged.

Every resulting \(J\)-order is well-ordered. Given a nonempty
subset \(X\) of \(J\), choose \(y\in X\) and then a good
\(d\) above \(y\). The initial segment below \(d\) is
isomorphic to a well-ordered \(T\)-carrier. Thus
\(X\cap\{j:j\prec d\}\) has a least element. It is least
in all of \(X\), since every point outside this initial
segment is at least \(d\). Empty \(J\) is harmless.

There is a model with \(J\) of type \(W\). Choose a cofinal
set of ordinals \(\xi<W\) realized by \(T\)-carriers, and
one witnessing model at each such \(\xi\). At all other
indices choose any one of these models. Realize this set-indexed
family in disjoint fibres over the ordinal \(W\), and set
\(H(\xi,-,-)\) to an isomorphism onto its initial segment
at each chosen \(\xi\). Interpret it as empty elsewhere.
Because \(W\) is a limit ordinal, the good indices are
strictly above every index. This satisfies the new theory and
pins \(W\), contrary to \cref{thm:p03-op-initial}.
Hence a theory with the asserted cofinal models must have a
non-well-ordered carrier model.

For every \(\gamma<W\), choose a pinnable \(\alpha\geq\gamma\).
The initial-segment construction in \cref{thm:p03-op-initial}
has only well-ordered models and realizes every type at most
\(\alpha\). Its types below \(\gamma\) are therefore cofinal
in \(\gamma\), so it refutes the threshold property there.
At \(\gamma=0\) the cofinality condition is empty and the same
theory refutes it. If the pinnable ordinals are unbounded,
this argument applies to every ordinal \(\gamma\).

For fields, use their definable carriers for the upper-bound
argument. For every infinite \(\gamma<W\), the
initial-segment construction followed by the field restriction
still gives types cofinal below \(\gamma\): omission of the
single type one does not affect cofinality below an infinite
ordinal. It therefore refutes every smaller infinite threshold.
The obstruction at finite \(\gamma=2\) explains why the
field version cannot assert a least ordinal without qualification.
\end{proof}
\end{annexitem}

\begin{annexitem}{annex:p03.op.3}{thm:p03-op-arithmetic}{Ordinal sum, product and exponentiation}
\begin{proof}
Place the two pinning theories in fibres at two distinguished indices,
using \cref{thm:p03-general-fibres} with disjoint copies of their
vocabularies. Their interpreted carriers are strict well-orders
\(A,<_{A}\) and \(B,<_{B}\), of possibly different types in
different models. Every pair of witnessing models can be realized
in this way, and the ambient domain may include an additional
set of coding points. The fibre projections send these extra
points to fixed default points, leaving both induced structures
unchanged. The two translated theories and finitely many
first-order clauses still have the required sentence bound,
by the finite-conjunction argument of \cref{thm:p03-op-initial}.
All constructions below concern the designated carriers, not
necessarily the whole fibre domains.

\emph{Sum.} Add disjoint carriers \(A',B'\) and graphs of
bijections from \(A,B\) to them. Order each by its transported
order, and put every point of \(A'\) before every point of
\(B'\). First-order biconditionals specify the resulting order
on \(A'\cup B'\). Every model has a well-order of the sum of
its two factor types. Disjoint copies of the witness carriers
realize type \(\alpha+\beta\).

\emph{Product.} Add a carrier \(C\) with two projection
functions. Require that they give a bijection from \(C\) to
\(B\times A\): each code has coordinates in the two carriers,
every pair has exactly one code, and equality of both coordinates
implies equality of codes. Define its order lexicographically,
first comparing the \(B\)-coordinate, then the \(A\)-coordinate.
These are finite first-order conditions. They give a well-order
of type \(\operatorname{otp}(A)\operatorname{otp}(B)\),
and the witness factors give \(\alpha\beta\). Empty factors
give empty \(C\), without making the ambient universe empty.

\emph{Exponentiation.} Suppose first that \(A\) is nonempty,
and let \(z\) be its least element. Add a carrier \(E\), a
point \(e\in E\), and functions \(q,b,a\). The intended
meaning of a nonempty code \(u\) is its prefix \(q(u)\),
last index \(b(u)\in B\), and nonzero coefficient
\(a(u)\in A\setminus\{z\}\). Impose the following finite
first-order conditions:
\begin{enumerate}
\item For \(u\in E\setminus\{e\}\), the three indicated
values belong to those carriers, \(q(u)\in E\), and
\(q(u)\ne e\) implies \(b(q(u))<_{B}b(u)\).
\item For each \(v\in E\), \(j\in B\), and
\(c\in A\setminus\{z\}\) such that
\(v=e\) or \(b(v)<_{B}j\), there is exactly one
\(u\in E\setminus\{e\}\) with
\(q(u)=v\), \(b(u)=j\), and \(a(u)=c\).
\end{enumerate}
Every nonempty code is among the codes in the second clause,
by the first. A sequence of successive prefixes that never
reaches \(e\) would give an infinite strictly decreasing
sequence of indices in the well-order \(B\). Therefore every
code reaches \(e\) after finitely many steps. Conversely,
the existence-and-uniqueness clause constructs exactly one
code for each finite list of increasing indices and nonzero
coefficients. Induction on list length proves both uniqueness
and that these are all codes. No axiom asserting finiteness of
codes has been assumed. It follows from well-foundedness of
the factor order.

Add a relation \(\lhd\) on \(E\). Require \(e\) to be
below every other code and no code to be below \(e\).
For \(u,v\ne e\), require
\[
\begin{aligned}
 u\lhd v\quad\Longleftrightarrow\quad& b(u)<_{B}b(v)\\
 &{}\lor\bigl(b(u)=b(v)\land a(u)<_{A}a(v)\bigr)\\
 &{}\lor\bigl(b(u)=b(v)\land a(u)=a(v)
                          \land q(u)\lhd q(v)\bigr).
\end{aligned}
\]
Require it to hold only on pairs from \(E\). This is a single
first-order recursive comparison clause, not an additional
recursion operation in the logic. Induction on the sum of the
two finite code lengths determines every comparison uniquely:
it compares finite-support functions from \(B\) to \(A\)
at the largest index where their values differ, with missing
values interpreted as \(z\).

Write \(\xi=\operatorname{otp}(A)>0\) and
\(\eta=\operatorname{otp}(B)\). This order of finite-support
functions has type \(\xi^\eta\), as follows by transfinite
induction on \(\eta\). For zero indices there is only the
empty function, giving type one. At a successor index, group
functions by the value at that last index. Each group has the
order for the earlier indices, and the groups occur in order
type \(\xi\), so the new type is
\(\xi^\gamma\xi=\xi^{\gamma+1}\).
At a limit, every finite support is bounded below that limit.
The earlier orders form nested initial segments: a function
with a nonzero coefficient beyond an earlier segment exceeds
every function supported inside it. Their union therefore has
type \(\sup_{\gamma<\eta}\xi^\gamma=\xi^\eta\).
This also covers \(\xi=1\), when only the empty code exists.
It proves well-ordering in every model and the claimed type
in a witness. In a witness, use the actual set of finite-support
functions as additional ambient codes, so the existence clauses
are satisfiable.

Finally allow empty \(A\) by a first-order case distinction:
if \(A\) is empty and \(B\) is empty, require the output
carrier to be a singleton; if \(A\) is empty and \(B\) is
nonempty, require it to be empty. Otherwise use the preceding
clauses with a least point \(z\). The point and auxiliary
functions are total on the ambient nonempty domain and are
unconstrained outside their guarded cases. This realizes
\(0^0=1\) and \(0^\eta=0\) for \(\eta>0\).
Finite Boolean combinations keep this one finite package of
first-order clauses. Thus the resulting theory pins
\(\alpha^\beta\) in all cases. The final assertion follows
from the exact description of pinnable ordinals below a bounded
supremum in \cref{thm:p03-op-initial}.
\end{proof}
\end{annexitem}
